%======================================================%
%   Beamer Presentation — C*-代数的表示
%   compile using PDFTeXify or PDFLaTeX
%======================================================%

%--------------------------------------------------------------------------------
%   PACKAGES AND THEMES
%--------------------------------------------------------------------------------

%\documentclass[notheorems,11pt,compress]{beamer}
\documentclass[notheorems,11pt]{beamer}

\usetheme{moloch}
\usepackage{appendixnumberbeamer}

\usepackage{booktabs}
\usepackage[scale=2]{ccicons}

\usepackage{pgfplots}
\usepgfplotslibrary{dateplot}

\usepackage[english]{babel}
\molochset{block=fill}
\usepackage{xspace}

\usepackage{physics2}
\usephysicsmodule{ab.legacy,op.legacy,diagmat}

\usefonttheme{professionalfonts}

\setbeamertemplate{blocks}[shadow=true]
\setbeamertemplate{navigation symbols}{}

\definecolor{iron}{RGB}{0,82,67}
\definecolor{footed}{rgb}{0.5,0.2,0.5}
\setbeamercolor{frametitle}{bg=iron}

%分式使用行间版
\everymath{\displaystyle}

\setbeamertemplate{footline}
{%
\leavevmode%
\hfill\normalfont\color{footed}\insertframenumber\,/\,\inserttotalframenumber \kern1em \hfill
\vskip4pt%
}


%----------- 宏包 -------------
\usepackage[UTF8,noindent]{ctex}
\usepackage{amsmath,amssymb,version}
\usepackage{graphicx,fancybox,mathrsfs,multirow}
\usepackage{booktabs}
\usepackage{epsfig,epstopdf}
\usepackage{url}
\usepackage{hyperref}
\hypersetup{
    colorlinks=true,
    linkcolor=iron,
    filecolor=magenta,
    urlcolor=cyan,
    citecolor=blue,
}
\usepackage{tabularx,array,makecell}
\usepackage{color,xcolor}
\usepackage{mathtools}
\usepackage{tikz}
\usepackage{makecell}

\usepackage{scrextend}

%---------- 定义表格命令 ----------
\newcolumntype{P}[1]{>{\centering \arraybackslash}p{#1}}
\newcolumntype{L}{X}
\newcolumntype{C}{>{\centering \arraybackslash}X}
\newcolumntype{R}{>{\raggedleft \arraybackslash}X}


%---------- 设置 定理环境  --------------
\theoremstyle{plain}
\setbeamertemplate{theorems}[numbered]
\newtheorem{theorem}{定理}
\numberwithin{theorem}{section}
\newtheorem{definition}{定义}
\newtheorem{lemma}{引理}
\numberwithin{lemma}{section}
\newtheorem{proposition}{命题}
\newtheorem{corollary}{{推论}}
\numberwithin{corollary}{section}
\theoremstyle{example}
\newtheorem{example}{例}
\renewenvironment{proof}[1][证明]{\textbf{#1}:~}{\qed\par}

\setbeamerfont{caption name}{family=\rmfamily,series=\normalfont}
\setbeamertemplate{caption}[numbered]

%---------- 定义行距 ----------
\renewcommand{\baselinestretch}{1.15}

\makeatletter
\newcommand\HUGE{\@setfontsize\Huge{36}{42}}
\makeatother

%================ 缩小公式上下间距 ================
\AtBeginDocument{
  \abovedisplayskip=5pt plus 1pt minus 1.5pt
  \belowdisplayskip=5pt plus 1pt minus 1.5pt
  \abovedisplayshortskip=0pt plus 3pt
  \belowdisplayshortskip=4pt plus 1pt minus 1.5pt
}

%---------- 自定义命令 ----------
\newcommand{\red}[1]{\textcolor{red}{#1}}
\newcommand{\blue}[1]{\textcolor{blue}{#1}}
\newcommand{\hl}[1]{\textcolor{orange}{#1}}

%---------- 算子 ----------
\DeclareMathOperator{\Spec}{Spec}


%----------------------------------------------------------------------------------------
%	TITLE PAGE
%----------------------------------------------------------------------------------------

\title[$C^*$代数的表示]{$C^*-\!$代数的表示：从 Wedderburn--Artin 到 Gelfand 变换}

\author{汇报人：庞逸天}
\institute[SJTU]
{
上海交通大学 \\
\medskip
\textit{yitianpang@sjtu.edu.cn}
}
\date[2026.6。15]{2026~年~6~月~15~日}

\graphicspath{{./figures/}}

\renewcommand{\familydefault}{\rmdefault}

\begin{document}

%------------------------------------------------
% TITLE
%------------------------------------------------

\begin{frame}
\titlepage
\end{frame}

\begin{frame}{目录}
  \setbeamertemplate{section in toc}[sections numbered]
\tableofcontents
\end{frame}

%=====================================================================
\section{引言}
%=====================================================================

\begin{frame}{群表示与代数表示}
  \begin{block}{课程主线}
    有限群表示 $\longrightarrow$ 一般代数上的模论 $\longrightarrow$ \hl{Wedderburn--Artin}：半单代数 $\cong \bigoplus M_{n_i}(\mathbb{C})$
  \end{block}
  \pause
  \begin{block}{一个自然的问题}
    如果额外赋予\alert{分析结构}（范数 + 对合），半单代数的分类会发生什么改变？
  \end{block}
  \pause
  \begin{itemize}
    \item 有限维：$C^*$条件是否改变分类？
    \item 交换：特征标群的无穷维推广 —— Gelfand 谱
  \end{itemize}
\end{frame}

%=====================================================================
\section{从矩阵代数到 $C^*-\!$代数}
%=====================================================================

\begin{frame}{矩阵代数的两个结构}
  $M_n(\mathbb{C})$ 除了矩阵乘法，还天然携带：
  \begin{enumerate}
    \item \textbf{对合} $A \mapsto A^*$（共轭转置）：$(AB)^* = B^*A^*,\;(A^*)^* = A$
    \item \textbf{算子范数} $\|A\| := \sup_{\|x\|=1} \|Ax\|$，后续选择谱范数$\|A\|_2=\max_{\lambda\in\sigma(A)}|\lambda|$
  \end{enumerate}
  \pause
  \begin{alertblock}{核心相容关系}
    \begin{equation*}
      \boxed{\|A^*A\| = \|A\|^2} \quad \forall A \in M_n(\mathbb{C})
    \end{equation*}
  \end{alertblock}
  将代数结构（对合）与分析结构（范数）紧密耦合在一起.
\end{frame}

\begin{frame}{$C^*$恒等式的两种证明}
  \textbf{证明一（泛函分析视角）}：$A^*A$ 自伴，故\par
  {\small$\|A^*A\| = \sup_{\|x\|=1}\langle A^*Ax,x\rangle =\sup_{\|x\|=1}\langle Ax,Ax\rangle= \sup_{\|x\|=1}\|Ax\|^2 = \|A\|^2$.}\par
  \pause
  \medskip
  \textbf{证明二（矩阵视角）}：考虑奇异值分解 $A = U\Sigma V^*$，其中 $\Sigma$ 是对角矩阵，$\sigma_1 \geq \sigma_2 \geq \cdots \geq 0$ 是 $A$ 的奇异值.
  则
  \begin{equation*}
    A^*A = V\Sigma^2 V^* \;\Rightarrow\; \|A^*A\| = \sigma_1^2 = \|A\|^2
  \end{equation*}
  \pause
  \begin{block}{抽象化}
    将这两个结构公理化 $\longrightarrow$ \alert{$C^*-\!$代数}
  \end{block}
\end{frame}

\begin{frame}{$C^*-\!$代数的定义}
  \begin{definition}[${}^*-\!$代数]
    复代数 $A$，配备对合 $a \mapsto a^*$ 满足：
    \begin{itemize}
      \item $(a+b)^* = a^*+b^*$，$(\lambda a)^* = \overline{\lambda}a^*$
      \item $(ab)^* = b^*a^*$（反同态），$(a^*)^* = a$
    \end{itemize}
  \end{definition}
  \pause
  \begin{definition}[$C^*-\!$代数]
    一个 ${}^*-\!$代数 $A$ 配备完备范数 $\|\cdot\|$ 满足：
    \begin{itemize}
      \item $\|ab\| \le \|a\|\|b\|$（次乘性）
      \item $\|a^*a\| = \|a\|^2$（\alert{$C^*$恒等式}）
    \end{itemize}
  \end{definition}
\end{frame}

\begin{frame}{关键性质：范数由对合唯一决定}
  \begin{definition}[谱]
    设 $A$ 有单位元 $1$.对 $a \in A$，定义其谱为
\begin{equation*}
    \sigma(a) := \{\lambda \in \mathbb{C} \mid a - \lambda \cdot 1 \text{ 在 } A \text{ 中不可逆}\}.
\end{equation*}
谱半径定义为 $\rho(a) := \sup\{|\lambda| : \lambda \in \sigma(a)\}$.
  \end{definition}
  \begin{proposition}
    $C^*-\!$代数中，$\|a\| = \sqrt{\rho(a^*a)}$.
  \end{proposition}
  \pause
  证明工具：\alert{Gelfand--Beurling 谱半径公式}
  \begin{equation*}
    \rho(a) = \lim_{n\to\infty} \|a^n\|^{1/n}
  \end{equation*}
\end{frame}

\begin{frame}{范数由对合唯一决定}
  \setcounter{proposition}{0}
  \begin{proposition}
    $C^*-\!$代数中，$\|a\| = \sqrt{\rho(a^*a)}$.
  \end{proposition}
  \begin{proof}
    对自伴元 $h$，由$C^*-\!$恒等式, $\|h^2\| = \|h^*h\| = \|h\|^2$，归纳可得 $\|h^{2^n}\| = \|h\|^{2^n}$.\par
    由 Gelfand--Beurling 谱半径公式 $\rho(h) = \lim_{k \to \infty} \|h^k\|^{1/k}$，取子列 $k = 2^n$，即得 $\rho(h) = \|h\|$.\par
    对一般 $a$，注意到 $a^*a$ 恒为自伴元，故
    \begin{equation*}
        \|a\| = \sqrt{\|a^*a\|} = \sqrt{\rho(a^*a)}.
    \end{equation*}
\end{proof}
  \begin{corollary}
    $C^*$代数上的范数被其 ${}^*-\!$代数结构\alert{唯一决定}.
  \end{corollary}
\end{frame}

\begin{frame}{对合是等距映射}
  \begin{proposition}\label{prop:isometry}
    在 $C^*$-代数中，$\|a^*\| = \|a\|$ 对所有 $a \in A$ 成立.
  \end{proposition}
  \pause
  \begin{proof}\small
    由 $C^*$恒等式和次乘性，$\|a\|^2 = \|a^*a\| \le \|a^*\|\|a\|$，故 $\|a\| \le \|a^*\|$.\par
    将 $a$ 替换为 $a^*$ 并利用 $(a^*)^* = a$，得 $\|a^*\| \le \|a\|$.两不等式合并即得等号.
  \end{proof}
\end{frame}



%=====================================================================
\section{有限维 $C^*-\!$代数的分类}
%=====================================================================

\begin{frame}{有限维 $C^*-\!$代数的半单性}
  \setcounter{lemma}{1}
  \begin{lemma}[半单性]\label{lem:semisimple-beamer}
    有限维 $C^*-\!$代数必为半单代数.
  \end{lemma}
  \pause
  \begin{proof}\small
    设 $J = \mathrm{rad}(A)$ 为 Jacobson 根.有限维代数的 Jacobson 根为幂零理想.
    若 $J \neq 0$，取 $a \in J \setminus \{0\}$，则 $a^*a \in J$ 幂零：存在 $m$ 使 $(a^*a)^m = 0$.\par
    取 $k$ 使 $2^k \ge m$，反复用 $C^*$恒等式：$\|a^*a\|^{2^k} = \|(a^*a)^{2^k}\| = 0$，故 $\|a^*a\| = 0$.\par
    由 $C^*$恒等式 $\|a\|^2 = \|a^*a\| = 0$，得 $\|a\| = 0$，矛盾.\par
    故 $J = 0$，$A$ 半单.
  \end{proof}
\end{frame}


\begin{frame}{分类定理}
  \begin{theorem}[有限维 $C^*-\!$代数的结构]
    设 $A$ 为有限维含幺 $C^*-\!$代数，则存在 $n_1,\dots,n_t \in \mathbb{N}^+$ 使得
    \begin{equation*}
      A \cong \bigoplus_{i=1}^t M_{n_i}(\mathbb{C})
    \end{equation*}
    为等距 ${}^*-\!$同构，且 $M_{n_i}(\mathbb{C})$ 上的对合为共轭转置.
  \end{theorem}
  \pause
  \begin{proof}\small
    \textbf{第1步.} 由半单性，Wedderburn--Artin 给出 $A \cong \bigoplus_i M_{n_i}(D_i)$，$D_i$ 为 $\mathbb{C}$ 上有限维可除代数.\par
    \textbf{第2步.} $\mathbb{C}$ 是代数闭域，其上有限维可除代数只有 $\mathbb{C}$ 自身.故 $D_i \cong \mathbb{C}$.\par
    \textbf{第3步.} 首先利用同构$\psi$构造 $M_n(\mathbb{C})$上的对合$y^\dag := \psi\bigl(\psi^{-1}(y)^*\bigr)$.可以证明 $M_n(\mathbb{C})$ 上 $C^*$对合在 ${}^*-\!$同构意义下唯一（共轭转置）；命题1证明范数由对合唯一决定.故同构自动保持对合与范数.
  \end{proof}
\end{frame}


%=====================================================================
\section{Gelfand 表示理论}
%=====================================================================

\begin{frame}{动机：特征标群的推广}
  \begin{block}{课程已知}
    有限交换群的不可约表示均为 1 维，构成特征标群 $\widehat{G}$.
  \end{block}
  \pause
  \begin{block}{$C^*-\!$代数的类比}
    交换 $C^*-\!$代数的所有不可约表示均为 1 维 $\Leftrightarrow$ 非零代数同态 $A \to \mathbb{C}$.
    所有这样的同态构成 \alert{Gelfand 谱} $\Spec(A)$.
  \end{block}
  \pause
    \begin{definition}[Gelfand 谱]
    $\Spec(A) := \{\varphi: A \to \mathbb{C} \mid \varphi \text{ 是非零代数同态}\}$.
  \end{definition}
\end{frame}

\begin{frame}{Gelfand 谱：极大理想空间}
  \begin{theorem}[谱与极大理想]
    设 $A$ 为含幺交换 Banach 代数，则存在双射
    \[
      \varphi:\Spec(A)\longrightarrow \operatorname{Max}(A),\qquad
      \varphi\longmapsto \ker\varphi.
    \]
  \end{theorem}
  \pause
  \begin{proof}\footnotesize
    \fbox{$\ker\varphi$是极大理想.} 若 $0\neq\varphi:A\to\mathbb C$，则 $\varphi(1)=1$，且
    \[
      A/\ker\varphi\cong \operatorname{Im}\varphi=\mathbb C.
    \]
    商代数是域，故 $\ker\varphi$ 为极大理想.\par
    \fbox{满射.} 设 $M$ 极大.先用 Neumann 级数知 $1$ 附近的元素均可逆，所以真理想闭包仍不含 $1$；极大性推出 $M=\overline M$，即 $M$ 闭.
    因而 $A/M$ 是 Banach 可除代数，由 Gelfand--Mazur 定理
    \[
      A/M\cong\mathbb C.
    \]
    复合商映射 $A\twoheadrightarrow A/M$ 得到 $\varphi\in\Spec(A)$，且 $\ker\varphi=M$.\par
    \fbox{单射.} 若$\ker\varphi_1=\ker\varphi_2$，则$\forall\, a\in A$，$a-\varphi_1(a)1\in\ker\varphi_1$.故$\varphi_2(a-\varphi_1(a)1)=\varphi_2(a)-\varphi_2(1)\varphi_1(a)=\varphi_2(a)-\varphi_1(a)=0$
  \end{proof}
\end{frame}



\begin{frame}{谱 = Gelfand 谱上的点赋值}
  \begin{theorem}
    设 $A$ 为含幺交换 $C^*-\!$代数，$a \in A$.则
    \begin{equation*}
      \boxed{\sigma(a) = \{\varphi(a) \mid \varphi \in \Spec(A)\}}.
    \end{equation*}
  \end{theorem}
  \pause
  也就是说，元素 $a$ 的谱等于函数
  \[
    \widehat a:\Spec(A)\to\mathbb C,\qquad
    \widehat a(\varphi)=\varphi(a)
  \]
  的值域.
\end{frame}

\begin{frame}{谱 = 点赋值：证明}
  \begin{proof}\small
    固定 $\lambda \in \mathbb{C}$.
    \begin{itemize}
      \item 按谱定义：$\lambda\in\sigma(a)$ 当且仅当 $a-\lambda1$ 不可逆.
      \item 在交换含幺代数中，不可逆等价于生成的主理想 $(a-\lambda1)$ 是真理想.
      \item 每个真理想包含于某个极大理想 $M$，而上页定理给出 $M=\ker\varphi$.
    \end{itemize}
    \begin{align*}
      \lambda \in \sigma(a) &\iff a-\lambda\cdot 1 \text{ 不可逆} \\
      &\iff a-\lambda\cdot 1 \in \text{某极大理想} \\
      &\iff \exists\,\varphi \in \Spec(A),\; \varphi(a-\lambda\cdot 1) = 0 \\
      &\iff \exists\,\varphi \in \Spec(A),\; \varphi(a) = \lambda.
    \end{align*}
    最后一行用到 $\varphi(1)=1$.
    因此谱正是 $\widehat a:\varphi\mapsto\varphi(a)$ 的值域.
  \end{proof}
\end{frame}

\begin{frame}{Gelfand 谱的紧 Hausdorff 性质}
  \begin{theorem}
    含幺交换 $C^*-\!$代数 $A$ 的 Gelfand 谱 $\Spec(A)$ 在弱*拓扑下是紧 Hausdorff 空间.
  \end{theorem}
  \pause
  \begin{proof}\small
    \fbox{Hausdorff 性.} 若 $\psi_1\neq\psi_2$，存在 $a\in A$ 使 $\psi_1(a)\neq\psi_2(a)$.
    弱*拓扑使每个求值映射 $\psi\mapsto\psi(a)$ 连续，故可由 $\mathbb C$ 中不交开集拉回分离两点.\par
    \fbox{紧性第一步.} 对 $\varphi\in\Spec(A)$，有 $\varphi(1)=1$.
    又由谱值域关系的简单方向 $\varphi(a)\in\sigma(a)$，得
    $|\varphi(a)|\le\rho(a)\le\|a\|$，故 $\|\varphi\|=1$.
    因此 $\Spec(A)\subseteq B_{A^*}$；由 Banach--Alaoglu 定理，$B_{A^*}$ 弱*紧.\par
    \fbox{紧性第二步.} 在 $B_{A^*}$ 中，
    \begin{equation*}
      \Spec(A)=
      \{\psi:\psi(1)=1\}
      \cap\bigcap_{a,b\in A}\{\psi:\psi(ab)=\psi(a)\psi(b)\}.
    \end{equation*}
    每个条件都是弱*闭条件，故 $\Spec(A)$ 是弱*闭子集，从而紧.
  \end{proof}
\end{frame}

\begin{frame}{Gelfand 变换的定义}
  \begin{definition}[Gelfand 变换]
    设 $A$ 为含幺交换 $C^*-\!$代数.对每个 $a\in A$，定义函数
    \[
      \widehat a:\Spec(A)\longrightarrow\mathbb C,\qquad
      \widehat a(\varphi)=\varphi(a).
    \]
    映射
    \[
      \Gamma:A\longrightarrow C(\Spec(A)),\qquad
      a\longmapsto\widehat a
    \]
    称为 $A$ 的 \alert{Gelfand 变换}.
  \end{definition}
  \pause
  \begin{block}{为什么是连续函数？}
    $\Spec(A)$ 赋予弱*拓扑，而弱*拓扑正是使所有求值映射
    \[
      \varphi\longmapsto \varphi(a)
    \]
    连续的拓扑.因此 $\widehat a\in C(\Spec(A))$.
  \end{block}
\end{frame}

\begin{frame}{Gelfand--Naimark 定理（交换版）}
  \begin{theorem}[Gelfand--Naimark，交换版]
    设 $A$ 为含幺交换 $C^*-\!$代数，则 Gelfand 变换
    \[
      \Gamma:A\longrightarrow C(\Spec(A)),\qquad
      a\longmapsto\widehat a
    \]
    是一个 \alert{等距 ${}^*-\!$同构}.
  \end{theorem}
  \pause
  \begin{block}{定理的含义}
    每个抽象元素 $a\in A$ 都可以无损地看成谱空间上的连续函数 $\widehat a$.
    代数运算、对合与范数全部被保留下来.
  \end{block}
\end{frame}

\begin{frame}{Gelfand--Naimark 定理：${}^*-\!$同态}
  \begin{proof}\small
    \textbf{1. 代数同态：}
    $\widehat{ab}(\varphi)=\varphi(ab)=\varphi(a)\varphi(b)=\widehat{a}(\varphi)\widehat{b}(\varphi)$，
    且 $\widehat{1}=1$，线性也逐点成立.

    \medskip
    \textbf{2. 保持对合：}
    先将 $a=x+iy$，其中 $x=x^*,y=y^*$.
    自伴元谱为实数，故 $\varphi(x),\varphi(y)\in\mathbb R$，于是
    \[
      \widehat{a^*}(\varphi)=\varphi(x-iy)
      =\overline{\varphi(x+iy)}
      =\overline{\widehat a(\varphi)}.
    \]
    因而 $\Gamma$ 是 ${}^*-\!$同态.
     \let\qed\relax
  \end{proof}
\end{frame}

\begin{frame}{Gelfand--Naimark 定理：等距与满射}
  \begin{proof}\small
    \textbf{3. 等距.}
    由于\begin{equation*}
    \widehat{a}(\varphi)|^2=\overline{\widehat{a}(\varphi)}\widehat{a}(\varphi)=\widehat{a^*}(\varphi)\widehat{a}(\varphi)=\widehat{a^*a}(\varphi).\
    \end{equation*}
    故对一般 $a$，由 $C^*$ 恒等式和 $a^*a$ 自伴：
    \[
      \|\widehat a\|_\infty^2
      =\|\widehat{a^*a}\|_\infty
      =\rho(a^*a)
      =\|a^*a\|
      =\|a\|^2.
    \]
    故 $\Gamma$ 等距，特别地单射，且像 $\widehat A$ 是闭子代数.

    \medskip
    \textbf{4. 满射.}
    $\widehat A$ 含常数函数、分离 $\Spec(A)$ 中不同点，并且对复共轭封闭.
    由复 Stone--Weierstrass 定理，$\widehat A$ 在 $C(\Spec(A))$ 中稠密.
    但 $\widehat A$ 已由等距性知为闭集，所以
    \[
      \widehat A=C(\Spec(A)).
    \]
  \end{proof}
\end{frame}

\begin{frame}{例子：无限维交换情形}
  \begin{exampleblock}{$A=C[0,1]$}
    取交换 $C^*-\!$代数 $A=C[0,1]$，配备一致范数，
    对合为逐点复共轭 $f^*(x)=\overline{f(x)}$.
  \end{exampleblock}
  \pause
  \textbf{首先，考察底空间 $\Spec(A)$.}
  \begin{itemize}
    \item 对每个实数 $x\in[0,1]$，点赋值泛函
      \[
        \varphi_x:C[0,1]\to\mathbb C,\qquad \varphi_x(f)=f(x)
      \]
      自然是代数同态.
    \item 由紧性可严格证明：所有代数同态均由点赋值给出.
    \item 映射 $x\mapsto\varphi_x$ 构成拓扑同胚，故
      \[
        \Spec(C[0,1])\cong[0,1].
      \]
  \end{itemize}
\end{frame}

\begin{frame}{例子：$C[0,1]$ 的 Gelfand 变换}
  \textbf{其次，考察元素的谱.}
  \begin{itemize}
    \item 任取 $f\in C[0,1]$，其 Gelfand 变换在点 $\varphi_x$ 处为
      \[
        \widehat f(\varphi_x)=\varphi_x(f)=f(x).
      \]
    \item 根据定理“谱 = Gelfand 谱上的点赋值”，
      \[
        \sigma(f)
        =\{\varphi_x(f):x\in[0,1]\}
        =f([0,1]).
      \]
  \end{itemize}
  \pause
  \begin{block}{核心含义}
    连续函数代数中的谱，就是连续函数的值域.
    抽象的“不可逆性”在这里变成了“函数是否取到某个值”.
  \end{block}
\end{frame}



\begin{frame}{例子：循环群 $C_n$ 的对照}
  \begin{columns}[T,onlytextwidth]
    \begin{column}{0.48\textwidth}
      \begin{block}{Wedderburn 解释}
        $\mathbb C[C_n]$ 是有限维半单交换代数，矩阵块只能是 $1\times1$：
        \[
          \mathbb C[C_n]\cong \mathbb C^{\oplus n}.
        \]
      \end{block}
    \end{column}
    \begin{column}{0.48\textwidth}
      \begin{block}{Gelfand 解释}
        由特征标理论可得，谱空间是
        \[
          \Spec(\mathbb C[C_n])\cong\widehat{C_n}\cong \{1, \dots, n\}.
        \]
        故
        \[
          \mathbb C[C_n]\cong C(\{1, \dots, n\}).
        \]
      \end{block}
    \end{column}
  \end{columns}
  \pause
  显然
  \[
    \mathbb C^{\oplus n}\cong C(\{1,\dots,n\})
  \]
  Wedderburn--Artin 定理和Gelfand表示理论在有限维交换情形下给出完全一致的分类.
\end{frame}

%=====================================================================
\section{总结}
%=====================================================================

\begin{frame}{有限维和交换情形的对照}
  \begin{table}[H]
    \centering
    \renewcommand{\arraystretch}{1.35}
    \caption{有限维与交换 $C^*-\!$代数表示论的对比}
    \label{tab:comparison}
    \resizebox{\linewidth}{!}{%
    \begin{tabular}{c|c|c}
        \hline
        & \textbf{有限维} & \textbf{交换} \\
        \hline
        核心定理 & Wedderburn--Artin & Gelfand--Naimark（交换版） \\
        \hline
        结构 & $\displaystyle\bigoplus_{i=1}^t M_{n_i}(\mathbb{C})$ & $C(\Spec(A))$ 或 $C_0(\Spec(A))$ \\
        \hline
        不可约表示 & $M_{n_i}(\mathbb{C})$ 上自由模 $\mathbb{C}^{n_i}$ & 1 维表示（点赋值 $\varphi_x$） \\
        \hline
        谱 $\sigma(a)$ & $a$ 的特征值集合 & $\widehat{a}$ 的值域 \\
        \hline
    \end{tabular}%
    }
\end{table}
\end{frame}


\begin{frame}[standout]
  \centering\LARGE
  谢谢！\\
  \smallskip
  欢迎提问
\end{frame}

\end{document}
